\section{Introduction}

Most users experience modern AI through a simple interface: a prompt goes in,
an answer comes back. Behind that simplicity, many of the most capable systems
are still delivered through a centralized service model. The 2026 AI Index
reports that industry produced more than 90\% of notable frontier models in
2025, and it also highlights how AI infrastructure remains concentrated around
large data centers and leading chip supply chains
\citep{stanford2026aiindex}. This model works well today, but it is not the
only possible shape for AI infrastructure. A plausible next step includes not
only remote frontier models, but also personal assistants, local models,
institutional models, domain-specific fine-tunes, and small services optimized
for narrow tasks. Recent work on collaborative inference, decentralized LLM
serving overlays, and personalized edge LLM agents points toward this more
diverse setting
\citep{borzunov2023petals,fang2026planetserve,wang2026ppai}.

If this capacity remains isolated, most users still depend on centralized
gateways. If it can be discovered, compared, routed to, and trusted,
individuals and organizations could share useful AI services while keeping
more control over where computation and data flow. The infrastructure question
is therefore not only \textit{which model is strongest?} It is also:
\textit{how should models be discovered, selected, trusted, and used without
overloading the same few providers?} LLM routing and profiling work already
treats model selection as a central problem
\citep{chen2023frugalgpt,ong2025routellm,xu2026routeprofile}.

This project, DIAL (Decentralized Intelligence Access Layer), studies one
concrete piece of that question: decentralized LLM service discovery and
routing. The core idea is to treat machines that host or access LLM backends
as network nodes. Each node publishes a capability profile, discovers other
peers through a peer-to-peer layer, and routes queries to nodes whose backend
is likely to answer well. A user can connect to any node and let the network
find a suitable executor without relying on a central model registry.

The prototype is intentionally practical. It focuses on a minimum useful
foundation: nodes join a local peer-to-peer network, publish capability
advertisements through a Kademlia DHT, exchange richer profiles, classify query
needs, score local and remote candidates, and answer locally or forward the
query. Around this layer, the project provides command-line tools, HTTP APIs,
DIAL Chat, and an evaluation harness. The central runtime path is therefore not
just ``send a prompt to another machine'': the entry node interprets the prompt
as capability demand, uses that demand to discover peers, compares full peer
profiles, and lets the selected node's own LLM backend produce the answer.
DIAL is therefore both a network prototype and an application-facing prototype:
the routing layer can be measured, and the same layer can be used through DIAL
Chat by a real user.

\begin{center}
\includegraphics[width=\textwidth]{dial_project_overview.png}
\captionof{figure}{Overview of DIAL. A user query reaches an entry node for this request, which classifies capability needs, uses DHT discovery and profile scoring to select a suitable peer, routes the query, and returns the answer generated by the selected node's LLM backend.}
\end{center}

\subsection{Project Thesis}

The main thesis is the following:

\begin{quote}
A peer-to-peer network is a practical starting point for decentralized LLM
service discovery and capability-aware query routing. DIAL focuses on the
first operational layer of that idea: finding peers, exchanging profiles, and
routing requests quickly. Blockchain was considered early in the project
because the goal was decentralized access, but the implemented design uses the
lighter P2P layer that better matches this first objective.
\end{quote}

This is why the prototype uses libp2p and a Kademlia DHT rather than a
blockchain-based design for the core network. Discovery is closer to BitTorrent
and IPFS than to payment settlement: nodes need fast, changing information
about who is available and what they can do
\citep{maymounkov2002kademlia,cohen2003bittorrent,benet2014ipfs}. Blockchain
may become useful later if the network needs a shared record of payments,
reputation, or other commitments between independent providers, but those
features come with extra cost and complexity and are not needed for the
discovery layer studied here \citep{nakamoto2008bitcoin,nassar2020blockchain}.

\subsection{Implemented Scope}

The project implements and evaluates the access layer of DIAL: decentralized
discovery, profile exchange, candidate scoring, and query forwarding. In
practice, this includes three main implementation outputs and one evaluation
workflow.

\begin{enumerate}
  \item A decentralized node architecture built around libp2p, a Kademlia
  DHT, profile exchange, capability scoring, and local or remote query
  execution.
  \item A capability-aware routing model that supports both single-capability
  and weighted multi-capability requests. It classifies a query into a small
  capability mix, uses that mix for DHT discovery, and scores candidates
  against the full mix while leaving room for future load-aware policies.
  \item An API layer and DIAL Chat web interface that make the network usable
  as an interactive application rather than only as a backend experiment.
  \item An evaluation workflow used to measure whether the routing layer behaves
  correctly under controlled local conditions.
\end{enumerate}

Stronger identity, privacy-preserving routing, abuse resistance, payments, and
reputation are important extensions, but they belong to later layers of the
system.

The implementation shows that a decentralized AI routing layer can be built
with current tools, while exposing the next problems: broader candidate
discovery, faster failed-peer handling, richer matching, answer aggregation,
incentives, privacy, and reputation.
